\section{AMC 8 2019}
        \begin{enumerate}
        	\item (\textit{AMC 8 2019}) Ike and Mike go into a sandwich shop with a total of $\$30.00$ to spend. Sandwiches cost $\$4.50$ each and soft drinks cost $\$1.00$ each. Ike and Mike plan to buy as many sandwiches as they can, and use any remaining money to buy soft drinks. Counting both sandwiches and soft drinks, how many items will they buy?


 $\textbf{(A) }6\qquad\textbf{(B) }7\qquad\textbf{(C) }8\qquad\textbf{(D) }9\qquad\textbf{(E) }10$

 \textbf{Answer: }D



\textbf{Solution 1}

We know that the sandwiches cost $4.50$ dollars. Guessing will bring us to multiplying $4.50$ by 6, which gives us $27.00$. Since they can spend $30.00$ they have $3$ dollars left. Since soft drinks cost $1.00$ dollar each, they can buy 3 soft drinks, which makes them spend $30.00$  Since they bought 6 sandwiches and 3 soft drinks, they bought a total of $9$ items. Therefore, the answer is $\boxed{\textbf{(D) }9}$.

 - SBose

 



\textbf{Solution 2 (Using Algebra)}

Let $s$ be the number of sandwiches and $d$ be the number of sodas. We have to satisfy the equation of
 \[4.50s+d=30\]In the question, it states that Ike and Mike buys as many sandwiches as possible. So, we drop the number of sodas for a while.We have: 
 \begin{align*} 4.50s&=30 \\ s&=\frac{30}{4.5} \\ s&=6R3 \end{align*}
We don't want a remainder so the maximum number of sandwiches is $6$.The total money spent is $6\cdot 4.50=27$.The number of dollar left to spent on sodas is $30-27=3$ dollars.$3$ dollars can buy $3$ sodas leading us to a total of $6+3=9$ items. Hence, the answer is $\boxed{\textbf{(D) }9}$.

 - by interactivemath

 



\textbf{Solution 3}

If the two could by fractions of sandwiches, they'd be able to buy $\frac{30.00\text{ dollars}}{4.50\frac{\text{dollars}}{\text{sandwich}}}=6\frac{2}{3}$ sandwiches. But they can only buy a whole number of sandwiches, so they can buy $6$, which costs $6\cdot\$4.50=\$27$. So they have $\$3$ left over, from which they can buy $3$ drinks (since they each cost $\$1$, for a total item count of $6+3=9$ items. The answer is $\boxed{\textbf{(D)}}$

 -Elihemi

 


	\item (\textit{AMC 8 2019}) Three identical rectangles are put together to form rectangle $ABCD$, as shown in the figure below. Given that the length of the shorter side of each of the rectangles is  5 feet, what is the area in square feet of rectangle $ABCD$?


 
                    \begin{figure}[H]
                        \centering
                        \includegraphics[width=0.5\textwidth]{Problem/AMC_8_AMC_2019_2_Problem_Figure_1.png}
                    \end{figure}
                    
$\textbf{(A) }45\qquad\textbf{(B) }75\qquad\textbf{(C) }100\qquad\textbf{(D) }125\qquad\textbf{(E) }150$

 \textbf{Answer: }E



\textbf{Solution 1}

We can see that there are $2$ rectangles lying on top of the other and that is the same as the length of one rectangle. Now we know that the shorter side is $5$, so the bigger side is $10$, if we do $5 \cdot 2 = 10$. Now we get the sides of the big rectangle being $15$ and $10$, so the area is $\boxed{\textbf{(E)}\ 150}$. ~avamarora

 



\textbf{Solution 2}

Using the diagram we find that the larger side of the small rectangle is $2$ times the length of the smaller side. Therefore, the longer side is $5 \cdot 2 = 10$. So the area of the identical rectangles is $5 \cdot 10 = 50$. We have 3 identical rectangles that form the large rectangle. Therefore the area of the large rectangle is $50 \cdot 3 = \boxed{\textbf{(E)}\ 150}$.    ~~fath2012

 



\textbf{Solution 3}

We see that if the short sides are 5, the long side has to be $5\cdot2=10$ because the long side is equal to the 2 short sides and because the rectangles are congruent. If that is to be, then the long side of the BIG rectangle(rectangle $ABCD$)is $10+5=15$ because long side + short side of the small rectangle is $15$. The short side of rectangle $ABCD$ is $10$ because it is the long side of the short rectangle. Multiplying $15$ and $10$ together gets us $15\cdot10$ which is $\boxed{\textbf{(E)}\ 150}$.~~mathboy282

 



\textbf{Solution 4}

We see that the $2$ rectangles lying on top of each other give us the height of the rectangle. Using what we know, we can find out that the $2$ rectangles put together is a square. So, we can infer that the length of the rectangles is $10$. Adding that to the width of the third rectangle which is $5$, we get that the length of the rectangle is $15$. Multiplying $10$ and $15$ gives us $15\cdot10$ which is $\boxed{\textbf{(E)}\ 150}$.~~awesomepony566

 



\textbf{Solution 5}

There are two rectangles lying on the side and one standing up. Given that one small side is 5, we can determine that two of the small sides make up a big side which means that the long side is equal to 10. The top side of the rectangle is made up of one small side and one long side, therefore the dimensions for the rectangle is 10x15. 10 multiplied by 15 is 150, hence the answer $\boxed{\textbf{(E)}\ 150}$.~elenafan

 $\textbf{Remark}$ALL THESE SOLUTIONS ARE VERY, VERY SIMILAR

 


	\item (\textit{AMC 8 2019}) Which of the following is the correct order of the fractions $\frac{15}{11},\frac{19}{15},$ and $\frac{17}{13},$ from least to greatest?


 $\textbf{(A) }\frac{15}{11}< \frac{17}{13}< \frac{19}{15}  \qquad\textbf{(B) }\frac{15}{11}< \frac{19}{15}<\frac{17}{13}    \qquad\textbf{(C) }\frac{17}{13}<\frac{19}{15}<\frac{15}{11}    \qquad\textbf{(D) } \frac{19}{15}<\frac{15}{11}<\frac{17}{13}   \qquad\textbf{(E) }   \frac{19}{15}<\frac{17}{13}<\frac{15}{11}$

 \textbf{Answer: }E



\textbf{Solution 1 (Bashing/Butterfly Method)}

We take a common denominator:
 \[\frac{15}{11},\frac{19}{15}, \frac{17}{13} = \frac{15\cdot 15 \cdot 13}{11\cdot 15 \cdot 13},\frac{19 \cdot 11 \cdot 13}{15\cdot 11 \cdot 13}, \frac{17 \cdot 11 \cdot 15}{13\cdot 11 \cdot 15} = \frac{2925}{2145},\frac{2717}{2145},\frac{2805}{2145}.\]

 Since $2717<2805<2925$ it follows that the answer is $\boxed{\textbf{(E)}\frac{19}{15}<\frac{17}{13}<\frac{15}{11}}$.

 Another approach to this problem is using the properties of one fraction being greater than another, also known as the \textbf{butterfly method}. That is, if $\frac{a}{b}>\frac{c}{d}$, then it must be true that $a * d$ is greater than $b * c$. Using this approach, we can check for \textbf{at least} two distinct pairs of fractions and find out the greater one of those two, logically giving us the expected answer of $\boxed{\textbf{(E)}\frac{19}{15}<\frac{17}{13}<\frac{15}{11}}$.

 -xMidnightFirex

 - Clearness by doulai1

 - A

 


	\item (\textit{AMC 8 2019}) Quadrilateral $ABCD$ is a rhombus with perimeter $52$ meters. The length of diagonal $\overline{AC}$ is $24$ meters. What is the area in square meters of rhombus $ABCD$?


 
                    \begin{figure}[H]
                        \centering
                        \includegraphics[width=0.5\textwidth]{Problem/AMC_8_AMC_2019_4_Problem_Figure_1.png}
                    \end{figure}
                    
$\textbf{(A) }60\qquad\textbf{(B) }90\qquad\textbf{(C) }105\qquad\textbf{(D) }120\qquad\textbf{(E) }144$

 \textbf{Answer: }D



\textbf{Solution 1}


                    \begin{figure}[H]
                        \centering
                        \includegraphics[width=0.5\textwidth]{Solution/AMC_8_AMC_2019_4_Solution_Figure_1.png}
                    \end{figure}
                    

 A rhombus has sides of equal length. Because the perimeter of the rhombus is $52$, each side is $\frac{52}{4}=13$. In a rhombus, diagonals are perpendicular and bisect each other, which means $\overline{AE}$ = $12$ = $\overline{EC}$.

 Consider one of the

 

                    \begin{figure}[H]
                        \centering
                        \includegraphics[width=0.5\textwidth]{Solution/AMC_8_AMC_2019_4_Solution_Figure_2.png}
                    \end{figure}
                    
$\overline{AB}$ = $13$, and $\overline{AE}$ = $12$. Using the Pythagorean theorem, we find that $\overline{BE}$ = $5$.You know the Pythagorean triple, (5, 12, 13).

 Thus the values of the two diagonals are $\overline{AC}$ = $24$ and $\overline{BD}$ = $10$.The area of a rhombus is = $\frac{d_1\cdot{d_2}}{2}$ = $\frac{24\cdot{10}}{2}$ = $120$
$\boxed{\textbf{(D)}\ 120}$




\textbf{Solution 2}

Right off the bat, we can see that the perimeter of the figure is 52. Dividing this by four, we can get that each side is equal to 13. By drawing a line perpendicular to the one given, we can split the figure into four right triangles. 12 (24/2) is equal to the height of one small right triangle, and 13 is the slanted side. Using the Pythagorean theorem we can find that 169 (13 squared) - 144 (12 squared) = 25 (five squared). With this, we can determine that each small right triangle equals 30. Multiplying that by four we can get $\boxed{\textbf{(D)}\ 120}$.

 


	\item (\textit{AMC 8 2019}) A tortoise challenges a hare to a race. The hare eagerly agrees and quickly runs ahead, leaving the slow-moving tortoise behind, believing that he will win, the hare stops to take a nap. Meanwhile, the tortoise walks at a slow, steady pace for the entire race. The hare wakes up and runs to the finish line, only to find the tortoise already there. Which of the following graphs matches the description of the race, showing the distance $d$ traveled by the two animals over time $t$ from start to finish?


 
                    \begin{figure}[H]
                        \centering
                        
\includegraphics[width=0.5\textwidth]{Problem/AMC_8_AMC_2019_5_Problem_Figure_1.png}
                    \end{figure}
                    

                    \begin{figure}[H]
                        \centering
                        
\includegraphics[width=0.5\textwidth]{Problem/AMC_8_AMC_2019_5_Problem_Figure_2.png}
                    \end{figure}
                    

 \textbf{Answer: }B



\textbf{Solution}

First, the tortoise walks at a constant rate, ruling out $(D)$. Second, when the hare is resting, the distance will stay the same, ruling out $(E)$ and $(C)$. Third, the tortoise wins the race, meaning that the non-constant one should go off the graph last, ruling out $(A)$. The answer $\boxed{\textbf{(B)}}$ is the only one left.

 


 


	\item (\textit{AMC 8 2019}) There are $81$ grid points (uniformly spaced) in the square shown in the diagram below, including the points on the edges. Point $P$ is in the center of the square. Given that point $Q$ is randomly chosen among the other $80$ points, what is the probability that the line $PQ$ is a line of symmetry for the square?


 
                    \begin{figure}[H]
                        \centering
                        \includegraphics[width=0.5\textwidth]{Problem/AMC_8_AMC_2019_6_Problem_Figure_1.png}
                    \end{figure}
                    
$\textbf{(A) }\frac{1}{5}\qquad\textbf{(B) }\frac{1}{4} \qquad\textbf{(C) }\frac{2}{5} \qquad\textbf{(D) }\frac{9}{20} \qquad\textbf{(E) }\frac{1}{2}$

 \textbf{Answer: }C



\textbf{Solution 1}


                    \begin{figure}[H]
                        \centering
                        \includegraphics[width=0.5\textwidth]{Solution/AMC_8_AMC_2019_6_Solution_Figure_1.png}
                    \end{figure}
                    Lines of symmetry go through point $P$, and there are $8$ directions the lines could go, and there are $4$ dots at each direction.$\frac{4\times8}{80}=\boxed{\textbf{(C)} \frac{2}{5}}$.

 



\textbf{Solution 2}

Divide the grid into 4 4x5 quadrants. Each row of 5 points has 1 point on a horizontal/vertical line of symmetry + 1 point on a diagonal line of symmetry: $\boxed{\textbf{(C)} \frac{2}{5}}$.

 


	\item (\textit{AMC 8 2019}) Shauna takes five tests, each worth a maximum of $100$ points. Her scores on the first three tests are $76$, $94$, and $87$. In order to average $81$ for all five tests, what is the lowest score she could earn on one of the other two tests?


 $\textbf{(A) }48\qquad\textbf{(B) }52\qquad\textbf{(C) }66\qquad\textbf{(D) }70\qquad\textbf{(E) }74$

 \textbf{Answer: }A



\textbf{Solution 1}

We should notice that we can turn the information we are given into a linear equation and just solve for our set variables. I'll use the variables $x$ and $y$ for the scores on the last two tests. 
 \[\frac{76+94+87+x+y}{5} = 81,\] 
 \[\frac{257+x+y}{5} = 81.\] We can now cross multiply to get rid of the denominator. 
 \[257+x+y = 405,\] 
 \[x+y = 148.\] Now that we have this equation, we will assign $y$ as the lowest score of the two other tests, and so: 
 \[x = 100,\] 
 \[y=48.\] Now we know that the lowest score on the two other tests is $\boxed{\textbf{(A)}\ 48}$.

 



\textbf{Solution 2}

Right now, she scored $76, 94,$ and $87$ points, for a total of $257$ points. She wants her average to be $81$ for her $5$ tests, so she needs to score $405$ points in total. This means she needs to score a total of $405-257=  148$ points in her next $2$ tests. Since the maximum score she can get on one of her $2$ tests is $100$, the least possible score she can get is $\boxed{\textbf{(A)}\ 48}$.

 Note: You can verify that $\boxed{48}$ is the right answer because it is the lowest answer out of the 5. Since it is possible to get 48, we are guaranteed that that is the right answer.

 



\textbf{Solution 3}

We can compare each of the scores with the average of $81$:$76$ $\rightarrow$ $-5$,$94$ $\rightarrow$ $+13$,$87$ $\rightarrow$ $+6$,$100$ $\rightarrow$ $+19$;

 So the last one has to be $-33$ (since all the differences have to sum to $0$), which corresponds to $81-33 = \boxed{48}$.

 



\textbf{Solution 4}

We know that she scored $76, 94,$ and $87$ points on her first $3$ tests for a total of $257$ points and that she wants her average to be $81$ for her $5$ total tests. Therefore, she needs to score a total of $405$ points. In addition, one of the final $2$ tests needs to be the maximum of $100$ points, to make the final test score—the one that we are looking for—the lowest score possible for her to earn. We can now see here that the sum of $76+94+100$ has a units digit of $0$ and that the final test score must have a units digit ending with a $5$. Now, $87$ needs to be added to a number that makes the sum divisible by $5$. Among the answer choices of $48, 52, 66, 70,$ and $74$, only $\boxed{\textbf{(A)}\ 48}$ has a units digit that works. ($48+87=135$, giving us a units digit of $5$.)

 



\textbf{Solution 5 (most intuitive)}

For Shauna to get an average of $81$, she needs to score a total of $81 \times 5 = 405$. One of the final two tests must earn a score of $100$ to ensure the final test score is the lowest possible. $76+94+87+100 = 357$, and $405-357 = \boxed{\textbf{(A)}\ 48}$.~themathkidwoohooooo

 



\textbf{Solution 6 (rocket speed)}

81 times 5 = 405. Adding the units digits of 76, 94, 87, and 100 gives 7. To end with five the fifth score ends with 8 so the answer is $\boxed{\textbf{(A)}\ 48}$. 

 


	\item (\textit{AMC 8 2019}) Gilda has a bag of marbles. She gives $20\%$ of them to her friend Pedro. Then Gilda gives $10\%$ of what is left to another friend, Ebony. Finally, Gilda gives $25\%$ of what is now left in the bag to her brother Jimmy. What percentage of her original bag of marbles does Gilda have left for herself?


 $\textbf{(A) }20\qquad\textbf{(B) }33\frac{1}{3}\qquad\textbf{(C) }38\qquad\textbf{(D) }45\qquad\textbf{(E) }54$

 \textbf{Answer: }E



\textbf{Solution 1}

After Gilda gives $20$\% of the marbles to Pedro, she has $80$\% of the marbles left. If she then gives $10$\% of what's left to Ebony, she has $(0.8*0.9)$ = $72$\% of what she had at the beginning. Finally, she gives $25$\% of what's left to her brother, so she has $(0.75*0.72)$ $\boxed{\textbf{(E)}\ 54}$ of what she had in the beginning left.

 



\textbf{Solution 2}

WOLG, suppose that Gilda has $100$ marbles. 

 Then, she gives Pedro 20\% of $100$, which equals $20$, leaving her with $80$ marbles.

 Out of $80$ marbles, she gives 10\% of $80$, which is 8 to Ebony. 

 Thus, she remains with $72$ marbles. 

 Then, she gives 25\% of $72$, which equals $18$ to Jimmy, finally leaving her with $54$. 

 And, $\frac{54}{100}$=54\%=$\boxed{\textbf{(E)}\ 54}$.

 



\textbf{Solution 3 (Cheese)}

Use this when you have about 10 minutes left and your brain does not function correctly.

 Since she gave away 20\% and 10\% of what is left and then another 25\% of what is actually left, we can do 20+10+25 or 55\%. But it is actually going to be a bit more than 55\% because 10\% of what is left is not 10\% of the total amount. So, the only option that is greater than 100\% - 55\% is $\boxed{\textbf{(E)}\ 54}$.

 


	\item (\textit{AMC 8 2019}) Alex and Felicia each have cats as pets. Alex buys cat food in cylindrical cans that are $6$ cm in diameter and $12$ cm high. Felicia buys cat food in cylindrical cans that are $12$ cm in diameter and $6$ cm high. What is the ratio of the volume of one of Alex's cans to the volume one of Felicia's cans?


 $\textbf{(A) }1:4\qquad\textbf{(B) }1:2\qquad\textbf{(C) }1:1\qquad\textbf{(D) }2:1\qquad\textbf{(E) }4:1$

 \textbf{Answer: }B



\textbf{Solution 1}

Using the formula for the volume of a cylinder $\pi r^2 h$, we get Alex, $108\pi$, and Felicia, $216\pi$. We can quickly notice that $\pi$ cancels out on both sides and that Alex's volume is $1/2$ of Felicia's leaving $1/2 = \boxed{\textbf{(B)}\ 1:2}$ as the answer. 

 



\textbf{Solution 2}

Using the formula for the volume of a cylinder, we get that the volume of Alex's can is $3^2\cdot12\cdot\pi$, and that the volume of Felicia's can is $6^2\cdot6\cdot\pi$. Now, we divide the volume of Alex's can by the volume of Felicia's can, so we get $\frac{1}{2}$, which is $\boxed{\textbf{(B)}\ 1:2}$.    

 -(Algebruh123)2020

 



\textbf{Solution 3}

The ratio of the numbers is $1/2$. Looking closely at the formula $r^2 * h * \pi$, we see that the $r * h * \pi$ will cancel, meaning that the ratio of them will be $\frac{1(2)}{2(2)}$ = $\boxed{\textbf{(B)}\ 1:2}$.

 -Lcz

 


 



\textbf{Solution 4}

The second can is $\cdot 2$ size in each of 2 dimensions, and $\cdot 1/2$ size in 1 dimension. $2^2/2 = \boxed{\textbf{(B)}\ 1:2}$.

 



\textbf{Solution 5}

Without calculating much, you can do($\pi ra^2) \cdot ha$ <-- which is Alex's volume, with ra being Alex's radius$(1/2 \cdot$ diameter), and $ha$ being her cylinders height$(\pi rf^2) \cdot hf <--$which is Felicia's volume, with $rf$ being Felicia's radius, and $hf$ being her cylinders height.Since we need the ratio between Alexa's and Felicias, we can do $(\pi ra^2)\cdot ha/(\pi rf^2)\cdot hf$ The $\pi$ cancel out, then substitute back in the numbers, which gives you:

 $(3^2 \cdot 12)/(6^2 \cdot 6) = (9 \cdot 12)/(36 \cdot 6) = 18/36 = 1/2 = 1:2$

 -wahahaqueenie

 


 


	\item (\textit{AMC 8 2019}) The diagram shows the number of students at soccer practice each weekday during last week. After computing the mean and median values, Coach discovers that there were actually $21$ participants on Wednesday. Which of the following statements describes the change in the mean and median after the correction is made?

                    \begin{figure}[H]
                        \centering
                        \includegraphics[width=0.5\textwidth]{Problem/AMC_8_AMC_2019_10_Problem_Figure_1.png}
                    \end{figure}
                    
$\textbf{(A) }$ The mean increases by $1$ and the median does not change.


 $\textbf{(B) }$ The mean increases by $1$ and the median increases by $1$.


 $\textbf{(C) }$ The mean increases by $1$ and the median increases by $5$.


 $\textbf{(D) }$ The mean increases by $5$ and the median increases by $1$.


 $\textbf{(E) }$ The mean increases by $5$ and the median increases by $5$.


 
 \textbf{Answer: }B



\textbf{Solution}

On Monday, $20$ people come. On Tuesday, $26$ people come. On Wednesday, $16$ people come. On Thursday, $22$ people come. Finally, on Friday, $16$ people come. $20+26+16+22+16=100$, so the mean is $20$. The median is $(16, 16, 20, 22, 26)$ $20$. The coach figures out that actually $21$ people come on Wednesday. The new mean is $21$, while the new median is $(16, 20, 21, 22, 26)$ $21$. Also, the median increases by $1$ because now the median is $21$ instead of $20$. The median and mean both change, so the answer is $\boxed{\textbf{(B)}}$.

 Another way to compute the change in mean is to notice that the sum increased by $5$ with the correction. So, the average increased by $5/5 = 1$. Then, the median is computed the same way.

 


	\item (\textit{AMC 8 2019}) The eighth grade class at Lincoln Middle School has $93$ students. Each student takes a math class, a foreign language class, or both. There are $70$ eighth graders taking a math class, and there are $54$ eighth graders taking a foreign language class. How many eighth graders take \textit{only} a math class and \textit{not} a foreign language class?


 $\textbf{(A) }16\qquad\textbf{(B) }23\qquad\textbf{(C) }31\qquad\textbf{(D) }39\qquad\textbf{(E) }70$

 \textbf{Answer: }D



\textbf{Solution 1}

Since each student must take a math class, a foreign language class, or both, all the students who do not take a foreign language class take a math class.  So $93 - 54 = \boxed{\textbf{(D)}\ 39}$ students take a math class but not a foreign language class. (The fact that $70$ students take a math class is irrelevant.)

 -hailstone, edited by j314andrews

 



\textbf{Solution 2}

Let $x$ be the number of students taking both a math and a foreign language class.

 By P-I-E, we get $70 + 54 - x$ = $93$. 

 Solving gives us $x = 31$.

 But we want the number of students taking only a math class,

 which is $70 - 31 = 39$.

 $\boxed{\textbf{(D)}\ 39}$

 
            \begin{enumerate}[label=(\roman*)]
                \item ):):):):):)
            \end{enumerate}
            



\textbf{Solution 3}

We have $70 + 54 = 124$ people taking classes. However, we over-counted the number of people who take both classes. If we subtract the original amount of people who take classes we get that $31$ people took the two classes. To find the amount of people who took only math class, we subtract the people who didn't take only one math class, so we get $70 - 31 = \boxed{\textbf{(D)} \, 39}$. 

 -fath2012

 



\textbf{Solution 4}


                    \begin{figure}[H]
                        \centering
                        \includegraphics[width=0.5\textwidth]{Solution/AMC_8_AMC_2019_11_Solution_Figure_1.png}
                    \end{figure}
                    

 We know that the sum of all three areas is $93$So, we have: 
 \[93 = 70-x+x+54-x\]
 \[93 = 70+54-x\]
 \[93 = 124 - x\]
 \[-31=-x\]
 \[x=31\]

 We are looking for the number of students in only math. This is $70-x$. Substituting $x$ with $31$, our answer is $\boxed{39}$.

 -mathnerdnair

 


	\item (\textit{AMC 8 2019}) The faces of a cube are painted in six different colors: red $(R)$, white $(W)$, green $(G)$, brown $(B)$, aqua $(A)$, and purple $(P)$. Three views of the cube are shown below. What is the color of the face opposite the aqua face?


 
                    \begin{figure}[H]
                        \centering
                        \includegraphics[width=0.5\textwidth]{Problem/AMC_8_AMC_2019_12_Problem_Figure_1.png}
                    \end{figure}
                    
$\textbf{(A) }\text{red}\qquad\textbf{(B) }\text{white}\qquad\textbf{(C) }\text{green}\qquad\textbf{(D) }\text{brown}\qquad\textbf{(E) }\text{purple}$

 \textbf{Answer: }A



\textbf{Solution 1}

$B$ is on the top, and $R$ is on the side, and $G$ is on the right side. That means that (image $2$) $W$ is on the left side. From the third image, you know that $P$ must be on the bottom since $G$ is sideways. That leaves us with the back, so the back must be $A$. The front is opposite of the back, so the answer is $\boxed{\textbf{(A)}\ R}$.

 



\textbf{Solution 2}

Looking closely, we can see that all faces are connected with $R$ except for $A$. Thus, the answer is $\boxed{\textbf{(A)}\ R}$.

 It is A, just draw it out!~phoenixfire

 



\textbf{Solution 3}

From pic 1 and 2, we know that the G's opposite is W. From pic 1 and 3, the B's opposite is P. So the $A$'s opposite is $\boxed{\textbf{(A)}\ R}$.

 


	\item (\textit{AMC 8 2019}) A \textit{palindrome} is a number that has the same value when read from left to right or from right to left. (For example, 12321 is a palindrome.) Let $N$ be the least three-digit integer which is not a palindrome but which is the sum of three distinct two-digit palindromes. What is the sum of the digits of $N$?


 $\textbf{(A) }2\qquad\textbf{(B) }3\qquad\textbf{(C) }4\qquad\textbf{(D) }5\qquad\textbf{(E) }6$

 \textbf{Answer: }A



\textbf{Solution 1}

Note that the only positive 2-digit palindromes are multiples of 11, namely $11, 22, \ldots, 99$. Since $N$ is the sum of 2-digit palindromes, $N$ is necessarily a multiple of 11. The smallest 3-digit multiple of 11 which is not a palindrome is 110, so $N=110$ is a candidate solution. We must check that 110 can be written as the sum of three distinct 2-digit palindromes; this suffices as $110=77+22+11$. Then, $N = 110$, and the sum of the digits of $N$ is $1+1+0 = \boxed{\textbf{(A) }2}$.

 
            \begin{itemize}
        \item There are other sets of 2-digit numbers that satisfy this rule. Some of them are $110 = 11+33+66$ and $110 = 22+33+55$
            \end{itemize}
            



\textbf{Solution 2 (variant of Solution 1)}

We already know that two-digit palindromes can only be two-digit multiples of 11; which are: $11, 22, 33, 44, 55, 66, 77, 88,$ and $99$. Since this is clear, we will need to find out the least multiple of 11 that is not a palindrome. Then, we start counting. $110 \ldots$ Aha! This multiple of 11, 110, not only isn’t a palindrome, but it also is the sum of three distinct two-digit palindromes, for example: 11 + 22 + 77, 22 + 33 + 55, and 44 + 11 + 55! The sum of $N$’s digits is  $1+1+0 = \boxed{\textbf{(A) }2}$.

 Thank you to the writer of Solution 1 for inspiring me to create this!

 EarthSaver 15:13, 11 June 2021 (EDT)

 



\textbf{Solution 3 (basically a version of the above solutions)}

As stated above, two-digit palindromes can only be two-digit multiples of 11. We can see that if we add anything that are multiples of 11 together, we will again get a multiple of 11. For instance, $11+22=33$. Since we know this fact and we are finding the smallest value possible, we can start with the first three-digit multiple of 11 which is $110$. Since this is not a palindrome and can be the sum of 3 two-digit palindromes (see above solutions for more details), $110$ fits the bill. We can see that the sum of $110$ 's digits is $1+1+0 = \boxed{\textbf{(A) }2}$.

 



\textbf{Solution 4 (FINALLY smth different lol)}

Two digit palindromes are two digit multiples of 11. When we take $33 \times 3$, we get 99, therefore 4 numbers that work for this would be $33, 33, 33$ and $11$. To convert this into three numbers, we can subtract $11$ from the first $33$, then add both $11$s to the last $33$ to get the set $22, 33$ and $55$. $22+33+55 = 110$ and $1+1+0 = \boxed{\textbf{(A) }2}$.

 


	\item (\textit{AMC 8 2019}) Isabella has $6$ coupons that can be redeemed for free ice cream cones at Pete's Sweet Treats. In order to make the coupons last, she decides that she will redeem one every $10$ days until she has used them all. She knows that Pete's is closed on Sundays, but as she circles the $6$ dates on her calendar, she realizes that no circled date falls on a Sunday. On what day of the week does Isabella redeem her first coupon?


 $\textbf{(A) }\text{Monday}\qquad\textbf{(B) }\text{Tuesday}\qquad\textbf{(C) }\text{Wednesday}\qquad\textbf{(D) }\text{Thursday}\qquad\textbf{(E) }\text{Friday}$

 \textbf{Answer: }C



\textbf{Solution 1 ( Guess and Check! )}

Let $\text{Day }1$ to $\text{Day\\ }2$ denote a day where one coupon is redeemed and the day when the second coupon is redeemed. 

 If she starts on a $\text{Monday}$, she redeems her next coupon on $\text{Thursday}$. 

 $\text{Thursday}$ to $\text{Sunday}$.

 Thus, $\textbf{(A)}\ \text{Monday}$ is incorrect.

 
If she starts on a $\text{Tuesday}$, she redeems her next coupon on $\text{Friday}$.

 $\text{Friday}$ to $\text{Monday}$.

 $\text{Monday}$ to $\text{Thursday}$.

 $\text{Thursday}$ to $\text{Sunday}$.

 Thus $\textbf{(B)}\ \text{Tuesday}$ is incorrect.

 
If she starts on a $\text{Wednesday}$, she redeems her next coupon on $\text{Saturday}$.

 $\text{Saturday}$ to $\text{Tuesday}$.

 $\text{Tuesday}$ to $\text{Friday}$.

 $\text{Friday}$ to $\text{Monday}$.

 $\text{Monday}$ to $\text{Thursday}$.

 And on $\text{Thursday}$, she redeems her last coupon. 

 
No Sunday occured; thus, $\boxed{\textbf{(C)}\ \text{Wednesday}}$ is correct. 

 
Checking for the other options, 

 
If she starts on a $\text{Thursday}$, she redeems her next coupon on $\text{Sunday}$.

 Thus, $\textbf{(D)}\ \text{Thursday}$ is incorrect.

 
If she starts on a $\text{Friday}$, she redeems her next coupon on $\text{Monday}$.

 $\text{Monday}$ to $\text{Thursday}$.

 $\text{Thursday}$ to $\text{Sunday}$.

 
Checking for the other options gave us negative results; thus, the answer is $\boxed{\textbf{(C)}\ \text{Wednesday}}$.

 



\textbf{Solution 2}

Let the days of the week be numbers$\pmod 7$. $3$ and $7$ are coprime, so continuously adding $3$ to a number$\pmod 7$ will cycle through all numbers from $0$ to $6$. If a string of 6 numbers in this cycle does not contain $0$, then if you minus 3 from the first number of this cycle, it will always be $0$. So, the answer is $\boxed{\textbf{(C) Wednesday}}$.

 



\textbf{Solution 3}

Start counting on Sunday, to maximize the number of 10-day jumps before reaching Sunday again. Add the 10 ($\equiv 3 \pmod 7$) days to reach the first coupon day on $\boxed{\textbf{(C)\ Wednesday}}$.

 Note: This only works when 7 and 3 are relatively prime. Otherwise, you’d prefer to start on a day whose distance from Sunday is relatively prime to the jump size.

 



\textbf{Solution 4}

Let Sunday be Day 0, Monday be Day 1, Tuesday be Day 2, and so forth. We see that Sundays fall on Day $n$, where n is a multiple of seven. If Isabella starts using her coupons on Monday (Day 1), she will fall on a Day that is a multiple of seven, a Sunday (her third coupon will be "used" on Day 21). Similarly, if she starts using her coupons on Tuesday (Day 2), Isabella will fall on a Day that is a multiple of seven (Day 42). Repeating this process, if she starts on Wednesday (Day 3), Isabella will first fall on a Day that is a multiple of seven, Day 63 (13, 23, 33, 43, 53 are not multiples of seven), but on her seventh coupon, of which she only has six. So, the answer is $\boxed{\textbf{(C)}\text{ Wednesday}}$.

 


 



\textbf{Solution 5}

There is a ten day gap, so every time you are 3 days ahead. For Sunday to be the 7th day you must be 6 days ahead, so 6x3 = 18. 18 days prior is the same as 4 days prior to Sunday which is $\boxed{\textbf{(C)}\text{ Wednesday}}$.

 


 


	\item (\textit{AMC 8 2019}) On a beach, $50$ people are wearing sunglasses and $35$ people are wearing caps. Some people are wearing both sunglasses and caps. If one of the people wearing a cap is selected at random, the probability that this person is also wearing sunglasses is $\frac{2}{5}$. If instead, someone wearing sunglasses is selected at random, what is the probability that this person is also wearing a cap?


 $\textbf{(A) }\frac{14}{85}\qquad\textbf{(B) }\frac{7}{25}\qquad\textbf{(C) }\frac{2}{5}\qquad\textbf{(D) }\frac{4}{7}\qquad\textbf{(E) }\frac{7}{10}$

 \textbf{Answer: }B



\textbf{Solution 1}

The number of people wearing caps and sunglasses is $\frac{2}{5}\cdot35=14$. So then, 14 people out of the 50 people wearing sunglasses also have caps.

 $\frac{14}{50}=\boxed{\textbf{(B)}\frac{7}{25}}$




\textbf{Solution 2 (extremely freaking overcomplicated)}

Let $A$ be the event that a randomly selected person is wearing sunglasses, and let $B$ be the event that a randomly selected person is wearing a cap. We can write $P(A \cap B)$ in two ways: $P(A)P(B|A)$ or $P(B)P(A|B)$. Suppose there are $t$ people in total. Then 
 \[P(A) = \frac{50}{t}\] and 
 \[P(B) = \frac{35}{t}.\] Additionally, we know that the probability that someone is wearing sunglasses given that they wear a cap is $\frac{2}{5}$, so $P(A|B) = \frac{2}{5}$. We let $P(B|A)$, which is the quantity we want to find, be equal to $x$. Substituting in, we get 
 \[\frac{50}{t} \cdot x = \frac{35}{t} \cdot \frac{2}{5}\]
 \[\implies 50x = 35 \cdot \frac{2}{5}\]
 \[\implies 50x = 14\]
 \[\implies x = \frac{14}{50}\]
 \[= \boxed{\textbf{(B)}~\frac{7}{25}}\]

 Note: This solution makes use of the dependent events probability formula, $P(A \cap B) = P(A)P(B|A)$, where $P(B|A)$ represents the probability that $B$ occurs given that $A$ has already occurred and $P(A \cup B)$ represents the probability of both $A$ and $B$ happening.

 


	\item (\textit{AMC 8 2019}) Qiang drives $15$ miles at an average speed of $30$ miles per hour. How many additional miles will he have to drive at $55$ miles per hour to average $50$ miles per hour for the entire trip?


 $\textbf{(A) }45\qquad\textbf{(B) }62\qquad\textbf{(C) }90\qquad\textbf{(D) }110\qquad\textbf{(E) }135$

 \textbf{Answer: }D



\textbf{Solution 1 (by looking at the answer choices)}

The only option that is easily divisible by $55$ is $110$, which gives 2 hours of travel. And, the formula is $\frac{15}{30} + \frac{110}{55} = \frac{5}{2}$.

 And, $\text{Average Speed}$ = $\frac{\text{Total Distance}}{\text{Total Time}}$.

 Thus, $\frac{125}{50} = \frac{5}{2}$.

 Both are equal and thus our answer is $\boxed{\textbf{(D)}\ 110}.$




\textbf{Solution 2}

To calculate the average speed, simply evaluate the total distance over the total time. Let the number of additional miles he has to drive be $x.$ Therefore, the total distance is $15+x$ and the total time (in hours) is 
 \[\frac{15}{30}+\frac{x}{55}=\frac{1}{2}+\frac{x}{55}.\] We can set up the following equation: 
 \[\frac{15+x}{\frac{1}{2}+\frac{x}{55}}=50.\] Simplifying the equation, we get 
 \[15+x=25+\frac{10x}{11}.\] Solving the equation yields $x=110,$ so our answer is $\boxed{\textbf{(D)}\ 110}$.

 



\textbf{Solution 3}

If he travels $15$ miles at a speed of $30$ miles per hour, he travels for 30 min. Average rate is total distance over total time so $(15+d)/(0.5 + t) = 50$, where d is the distance left to travel and t is the time to travel that distance. Solve for $d$ to get $d = 10+50t$. You also know that he has to travel $55$ miles per hour for some time, so $d=55t$. Plug that in for d to get $55t = 10+50t$ and $t=2$ and since $d=55t$, $d = 2\cdot55 =110$, the answer is $\boxed{\textbf{(D)}\ 110}$.

 



\textbf{Solution 4}

Let $h$ be the amount of hours Qiang drives after his first 15 miles. Average speed, which we know is $50$ mph, means total distance over total time. For 15 miles at 30 mph, the time taken is $\frac{1}{2}$ hour, so the total time for this trip would be $\frac{1}{2} + h$ hours. For the total distance, 15 miles are traveled in the first part and $55h$ miles in the second. This gives the following equation: 

 

 \[\dfrac{15+55h}{\frac{1}{2}+h} = 50.\] 

 


 
Cross multiplying, we get that $15 + 55h = 50h + 25$, and simple algebra gives $h=2$. In 2 hours traveling at 55 mph, the distance traveled is $\frac{2 \hspace{0.05 in} \text{hours}}{1} \cdot \frac{55 \hspace{0.05 in} \text{miles}}{1 \hspace{0.05 in} \text{hour}} = 2 \cdot 55 \hspace{0.05 in} \text{miles} = 110 \hspace{0.05 in} \text{miles}$, which is choice $\boxed{\textbf{(D)}\ 110}$. 

 


 


	\item (\textit{AMC 8 2019}) What is the value of the product 


 
 \[\left(\frac{1\cdot3}{2\cdot2}\right)\left(\frac{2\cdot4}{3\cdot3}\right)\left(\frac{3\cdot5}{4\cdot4}\right)\cdots\left(\frac{97\cdot99}{98\cdot98}\right)\left(\frac{98\cdot100}{99\cdot99}\right)?\]
$\textbf{(A) }\frac{1}{2}\qquad\textbf{(B) }\frac{50}{99}\qquad\textbf{(C) }\frac{9800}{9801}\qquad\textbf{(D) }\frac{100}{99}\qquad\textbf{(E) }50$

 \textbf{Answer: }B



\textbf{Solution 1 (telescoping)}

We rewrite: 
 \[\frac{1}{2}\cdot\left(\frac{3\cdot2}{2\cdot3}\right)\left(\frac{4\cdot3}{3\cdot4}\right)\cdots\left(\frac{99\cdot98}{98\cdot99}\right)\cdot\frac{100}{99}\]

 The middle terms cancel, leaving us with

 
 \[\left(\frac{1\cdot100}{2\cdot99}\right)= \boxed{\textbf{(B)}\frac{50}{99}}\]




\textbf{Solution 2}

If you calculate the first few values of the equation, all of the values tend to be around $\frac{1}{2}$, but are not equal to it. The answer closest to $\frac{1}{2}$ but not equal to it is $\boxed{\textbf{(B)}\frac{50}{99}}$.

 



\textbf{Solution 3}

Rewriting the numerator and the denominator, we get $\frac{\frac{100! \cdot 98!}{2}}{\left(99!\right)^2}$. We can simplify by canceling 99! on both sides, leaving us with: $\frac{100 \cdot 98!}{2 \cdot 99!}$ We rewrite $99!$ as $99 \cdot 98!$ and cancel $98!$, which gets $\boxed{\textbf{(B)}\frac{50}{99}}$.

 



\textbf{Solution 4}

All of the terms have the form $\frac{k^2-1}{k^2}$, which is $<1$, so the product is $<1$, so we eliminate options (D) and (E). (C) is too close to 1 to be possible. The partial products seem to be approaching 1/2, so we guess that 1/2 is the limit/asymptote, and so any finite product would be slightly larger than 1/2. Therefore, by process of elimination and a small guess, we get that the answer is $\boxed{\textbf{(B)}\frac{50}{99}}$.

 



\textbf{Solution 5}

The product
 \[\left(\frac{1 \cdot 3}{2 \cdot 2}\right) \left(\frac{2 \cdot 4}{3 \cdot 3}\right) \left(\frac{3 \cdot 5}{4 \cdot 4}\right) \cdots \left(\frac{98 \cdot 100}{99 \cdot 99}\right)\]can be simplified by observing that in each individual fraction, the numerator and denominator contain factors that cancel out with adjacent terms. Specifically, the factor 3 in the numerator of the first fraction cancels with the 3 in the denominator of the second fraction, the 4 in the numerator of the second fraction cancels with the 4 in the denominator of the third fraction, and so on. This telescoping cancellation continues throughout the entire product.

 After all the cancellations, only the first factor of the first fraction, \( \frac{1}{2} \), and the last factor of the last fraction, \( \frac{100}{99} \), remain. The value of the product is therefore:
 \[\frac{1}{2} \cdot \frac{100}{99} = \frac{50}{99}.\]Thus, the value of the product is:(B) \(\frac{50}{99}\).

 -- Rayansh Mankad(SharpWhiz17)

 


 


	\item (\textit{AMC 8 2019}) The faces of each of two fair dice are numbered 1, 2, 3, 5, 7, and 8. When the two dice are tossed, what is the probability that their sum will be an even number?


 $\textbf{(A) }\frac{4}{9}\qquad\textbf{(B) }\frac{1}{2}\qquad\textbf{(C) }\frac{5}{9}\qquad\textbf{(D) }\frac{3}{5}\qquad\textbf{(E) }\frac{2}{3}$

 \textbf{Answer: }C



\textbf{Solution 1 (Casework)}

We have $2$ dice with $2$ evens and $4$ odds on each die. For the sum to be even, the 2 rolls must be $2$ odds or $2$ evens. 

 Ways to roll $2$ odds (Case $1$): The total number of ways to obtain $2$ odds on 2 rolls is $4*4=16$, as there are $4$ possible odds on the first roll and $4$ possible odds on the second roll.

 Ways to roll $2$ evens (Case $2$): Similarly, we have $2*2=4$ ways to obtain 2 evens. Probability is $\frac{20}{36}=\frac{5}{9}$, or $\framebox{C}$.

 



\textbf{Solution 2 (Complementary Counting)}

We count the ways to get an odd. If the sum is odd, then we must have an even and an odd. The probability of an even is $\frac{1}{3}$ , and the probability of an odd is $\frac{2}{3}$. We have to multiply by $2$ because the even and odd can be in any order. This gets us $\frac{4}{9}$, so the answer is $1 - \frac{4}{9} = \frac{5}{9} = \boxed{(\textbf{C})}$.

 



\textbf{Solution 3}

To get an even, you must get either 2 odds or 2 evens. The probability of getting 2 odds is $\frac{4}{6} \times \frac{4}{6}$. The probability of getting 2 evens is $\frac{2}{6} \times \frac{2}{6}$. If you add them together, you get $\frac{16}{36} + \frac{4}{36}$ = $\boxed{(\textbf{C})  \frac{5}{9}}$.

 


	\item (\textit{AMC 8 2019}) In a tournament there are six teams that play each other twice. A team earns $3$ points for a win, $1$ point for a draw, and $0$ points for a loss. After all the games have been played it turns out that the top three teams earned the same number of total points. What is the greatest possible number of total points for each of the top three teams?


 $\textbf{(A) }22\qquad\textbf{(B) }23\qquad\textbf{(C) }24\qquad\textbf{(D) }26\qquad\textbf{(E) }30$

 \textbf{Answer: }C



\textbf{Solution 1}

The teams play a total of $2 \cdot \binom{6}{2} = 30$ games, and $2 \cdot \binom{3}{2} = 6$ of these games only involve bottom-$3$ teams.  So $24$ games involve at least one top-$3$ team.  

 Since at most $3$ points are awarded in each game, the top-$3$ teams could earn a total of at most $24 \cdot 3 = 72$ points.  Thus it is impossible for the top $3$ teams to score more than $\frac{72}{3} = 24$ points each.

 It is possible for the top $3$ teams to score $\boxed{\textbf {(C) }24}$ points each.  This occurs if each top-$3$ team wins all its games against bottom-$3$ teams, and each top-$3$ team wins $1$ game and loses $1$ game against each of the other $2$ top-$3$ teams.

 



\textbf{Solution 2}

We can name the top three teams as $A$, $B$, and $C$. We can see that (respective scores of) $A=B=C$ because these teams have the same points. If we look at the matches that involve the top three teams, we see that there are some duplicates: $AB$, $BC$, and $AC$ come twice. In order to even out the scores and get the maximum score, we can say that in match $AB$, $A$ and $B$ each win once out of the two games that they play. We can say the same thing for $AC$ and $BC$. This tells us that each team $A$, $B$, and $C$ win and lose twice. This gives each team a total of $3 + 3 + 0 + 0 = 6$ points. Now, we need to include the other three teams. We can label these teams as $D$, $E$, and $F$. We can write down every match that $A, B,$ or $C$ plays in that we haven't counted yet: $AD$, $AD$, $AE$, $AE$, $AF$, $AF$, $BD$, $BD$, $BE$, $BE$, $BF$, $BF$, $CD$, $CD$, $CE$, $CE$, $CF$, and $CF$. We can say $A$, $B$, and $C$ win each of these in order to obtain the maximum score that $A$, $B$, and $C$ can have. If $A$, $B$, and $C$ win all six of their matches, $A$, $B$, and $C$ will have a score of $18$. $18 + 6$ results in a maximum score of $\boxed{\textbf{(C) }24}$




\textbf{Solution 3}

To start, we calculate how many games each team plays. Each team can play against $5$ other teams twice, so there are $10$ games that each team plays. So the answer is $10\cdot 3$ which is $30$ But wait... if we want $3$ teams to have the same amount of points, there can't possibly be a team who wins all their games. Let the top three teams be $A,B$, and $C.$ $A$ plays $B$ and $C$ twice, so in order to maximize the games being played, we can split it $50-50$ between the $4$ games $A$ plays against $B$ or $C$. We find that we just subtract $2$ games or $6$ points. Therefore the answer is $30-6$, $24$or $\boxed{\textbf{(C) }24}$




\textbf{Solution 4}

Let the top $3$ teams be $A, B, C$ and the bottom teams to be $D, E, F$. There are $2$ types of games the winners ($A, B, C$) can play: Winners vs losers and Winner vs Winner.

 Case 1: Winner vs Loser

 For A, it would be $AD, AE, AF$. We want to maximize this or make $A$ win in all $6$ ($3\times2$) games. Therefore, A would get $6\times3$ points or $18$.

 Case 2: Winner vs Winner

 There are $_{3}P_{2}=\frac{3!}{1!}=6$ games they can play against each other. Since we want each $3$ to have an equal amount. Each team should win $2$ games and get $6$ points. In total, $18+6=\boxed{\textbf{(C) }24}$~RandomMathGuy500

 



\textbf{Solution 5}

Name the teams A, B, C, D, E, and F.As we know that the top three teams, A, B, and C, tied, we can just focus on one, A. When A plays B or C, we want them to both win one and lose one, maximising the points for both teams. When A plays D, E, or F, we want A to win both times. This adds up to $3+3$(top teams)$+6+6+6$ (other teams), which equals 24 points.~ParticlePhysicsThis solution is by far the fastest and probably most used

 


	\item (\textit{AMC 8 2019}) How many different real numbers $x$ satisfy the equation 
 \[(x^{2}-5)^{2}=16?\]
$\textbf{(A) }0\qquad\textbf{(B) }1\qquad\textbf{(C) }2\qquad\textbf{(D) }4\qquad\textbf{(E) }8$

 \textbf{Answer: }D



\textbf{Solution 1}

We have that $(x^2-5)^2 = 16$ if and only if $x^2-5 = \pm 4$. If $x^2-5 = 4$, then $x^2 = 9 \implies x = \pm 3$, giving 2 solutions. If $x^2-5 = -4$, then $x^2 = 1 \implies x = \pm 1$, giving 2 more solutions. All four of these solutions work, so the answer is $\boxed{\textbf{(D) }4}$. Further, the equation is a \href{https://artofproblemsolving.com/wiki/index.php?title=Quartic\_Equation}{quartic} in $x$, so by the \href{https://artofproblemsolving.com/wiki/index.php/Fundamental\_Theorem\_of\_Algebra}{Fundamental Theorem of Algebra}, there can be at most four real solutions.

 



\textbf{Solution 2}

We can expand $(x^2-5)^2$ to get $x^4-10x^2+25$, so now our equation is $x^4-10x^2+25=16$. Subtracting $16$ from both sides gives us $x^4-10x^2+9=0$. Now, we can factor the left-hand side to get $(x^2-9)(x^2-1)=0$. If $x^2-9$ and/or $x^2-1$ equals $0$, then the entire left side will equal $0$. Since the solutions can be both positive and negative, we have $4$ solutions: $-3,3,-1,1$ (we can find these solutions by setting $x^2-9$ and $x^2-1$ equal to $0$ and solving for $x$). So, the answer is $\boxed{\textbf{(D) }4}$. Note: every positive root can also be shown as negative 

 



\textbf{Solution 3}

Subtract 16 from both sides and factor using difference of squares: 

 


 
 \[(x^2 - 5)^2 = 16\]
 \[(x^2 - 5)^2 - 16 =0\]
 \[(x^2 - 5)^2 - 4^2 = 0\] 
 \[[(x^2 - 5)-4][(x^2 - 5) + 4] = 0\]
 \[(x^2 - 9)(x^2 - 1) =0\]
 \[(x+3)(x-3)(x+1)(x-1) = 0.\]

 
Quite obviously, this equation has $\boxed{\textbf{(D) }4}$ solutions. 

 


	\item (\textit{AMC 8 2019}) What is the area of the triangle formed by the lines $y=5$, $y=1+x$, and $y=1-x$?


 $\textbf{(A) }4\qquad\textbf{(B) }8\qquad\textbf{(C) }10\qquad\textbf{(D) }12\qquad\textbf{(E) }16$

 \textbf{Answer: }E



\textbf{Solution 1}

First, we need to find the coordinates where the graphs intersect. 

 We want the points x and y to be the same. Thus, we set $5=x+1,$ and get $x=4.$ Plugging this into the equation, $y=1-x,$ $y=5$, and $y=1+x$ intersect at $(4,5)$, we call this line x.

 Doing the same thing, we get $x=-4.$ Thus, $y=5$. Also,$y=5$ and $y=1-x$ intersect at $(-4,5)$, and we call this line y.

 It's apparent the only solution to $1-x=1+x$ is $0.$ Thus, $y=1.$$y=1-x$ and $y=1+x$ intersect at $(0,1)$, we call this line z.

 Using the \href{https://artofproblemsolving.com/wiki/index.php?title=Shoelace\_Theorem}{Shoelace Theorem} we get: 
 \[\left(\frac{(20-4)-(-20+4)}{2}\right)=\frac{32}{2}\] $=$ So, our answer is $\boxed{\textbf{(E)}\ 16.}$

 We might also see that the lines $y$ and $x$ are mirror images of each other. This is because, when rewritten, their slopes can be multiplied by $-1$ to get the other. As the base is horizontal, this is an isosceles triangle with base 8, as the intersection points have a distance of 8. The height is $5-1=4,$ so $\frac{4\cdot 8}{2} = \boxed{\textbf{(E)} 16.}$

 Warning: Do not use the distance formula for the base then use Heron's formula. It will take you half of the time you have left or maybe more!!

 



\textbf{Solution 2}

Graphing the lines, using the intersection points we found in Solution 1, we can see that the height of the triangle is 4, and the base is 8. Using the formula for the area of a triangle, we get $\frac{4\cdot8}{2}$, which is equal to $\boxed{\textbf{(E)}\ 16}$.

 



\textbf{Solution 3}

$y = x + 1$ and $y = -x + 1$ have $y$-intercepts at $(0, 1)$ and slopes of $1$ and $-1$, respectively. Since the product of these slopes is $-1$, the two lines are perpendicular. From $y = 5$, we see that $(-4, 5)$ and $(4, 5)$ are the other two intersection points, and they are $8$ units apart. By symmetry, this triangle is a $45-45-90$ triangle, so the legs are $4\sqrt{2}$ each and the area is $\frac{(4\sqrt{2})^2}{2} = \boxed{\textbf{(E)}\ 16}$.

 


	\item (\textit{AMC 8 2019}) A store increased the original price of a shirt by a certain percent and then decreased the new price by the same percent. Given that the resulting price was $84\%$ of the original price, by what percent was the price increased and decreased?


 $\textbf{(A) }16\qquad\textbf{(B) }20\qquad\textbf{(C) }28\qquad\textbf{(D) }36\qquad\textbf{(E) }40$

 \textbf{Answer: }E



\textbf{Solution 1}

Suppose the fraction of discount is $x$. That means $(1-x)(1+x)=0.84$; so, $1-x^{2}=0.84$, and $(x^{2})=0.16$, procuring $x=0.4$. Therefore, the price was increased and decreased by $40$\%, or $\boxed{\textbf{(E)}\ 40}$.

 



\textbf{Solution 1.1}

After the first increase by $p$ percent, the shirt price became $(1+p)$ times greater than the original. Upon the decrease in p percent on this price, the shirt price became $(1-p)$ times less than $(1+p)$, or $(1-p)(1+p)$. We know that this price is $84$ percent of the original, so $(1-p)(1+p) = 0.84$. 

 


 From here, we can list the factors of $0.84$ and see which are equidistant from $1$. We see that $0.6$ and $1.4$ are both $0.4$ from $1$, so $p = 0.4 = 40 \%$, or choice $\boxed{\textbf{(E)}\ 40}$. 

 



\textbf{Solution 2 (Time-Consuming)}

We can try out every option and see which one works. Wasting 30 mins of the 40 mins, by this method, we get $\boxed{\textbf{(E)}\ 40}$.

 



\textbf{Solution 3}

Let x be the discount. We can also work in reverse such as ($84$)$(\frac{100}{100-x})$$(\frac{100}{100+x})$ = $100$.

 Thus, $8400$ = $(100+x)(100-x)$. Solving for $x$ gives us $x = 40, -40$. But $x$ has to be positive. Thus, $x$ = $40$.

 



\textbf{Solution 4 ~ using the answer choices (similar to Solution 2)}

Let our original cost be $\$ 100.$ We are looking for a result of $\$ 84,$ then. We try 16\% and see it gets us higher than 84. We try 20\% and see it gets us lower than 16 but still higher than 84. We know that the higher the percent, the less the value. We try 36, as we are not progressing much, and we are close! We try $\boxed{40\%}$, and we have the answer; it worked.(OR: try (C) first to eliminate 2 answer choices)

 



\textbf{Solution 5 (A Variation of Solution 4)}

Let our original cost be $\$ 100$, so we are looking for a whole number of $\$ 84$. Also, we can see that (A), (C), and (D) give us answers with decimals while we know that (B) and (E) give us whole numbers. Therefore, we only need to try these two: (B) $\$100$ increased by 20% = $\$120$, and $\$120$ decreased by 20% = $\$96$, a whole number, and (E) $\$100$ increased by 40% = $\$140$, and $\$140$ decreased by 40% = $\$84$, a whole number.

 Thus, $40$\% or $\boxed{\textbf{(E)}\ 40}$ is the answer.

 



\textbf{Solution 6 (using factors)}

First find the prime factors of 84. These are 2,2,3 and 7. Using this, try the answers. For example, 100*1.2/1.2. This is equal to 120*0.8. If you find the factors of 12 and 8, you find that there is no seven and too many 2's. So $\boxed{\textbf{(B)}\ 20}$ is out. Using the same method, you can find that 116*0.84 is out, as 84 already has all the factors of 84. 128*0.72 and 136*0.64 both don't have seven as a factor, so $\boxed{\textbf{(A)}\ 16}$, $\boxed{\textbf{(C)}\ 28}$, and $\boxed{\textbf{(D)}\ 36}$ is out.  This leaves 140*0.6, which has just the right amount of factors (2\textasciicircum{}2,3, and 7). Therefore, the answer is $\boxed{\textbf{(E)}\ 40}$



	\item (\textit{AMC 8 2019}) After Euclid High School's last basketball game, it was determined that $\frac{1}{4}$ of the team's points were scored by Alexa and $\frac{2}{7}$ were scored by Brittany. Chelsea scored $15$ points. None of the other $7$ team members scored more than $2$ points. What was the total number of points scored by the other $7$ team members?


 $\textbf{(A) }10\qquad\textbf{(B) }11\qquad\textbf{(C) }12\qquad\textbf{(D) }13\qquad\textbf{(E) }14$

 \textbf{Answer: }B



\textbf{Solution 1}

Given the information above, we start with the equation $\frac{t}{4}+\frac{2t}{7} + 15 + x = t$, where $t$ is the total number of points scored and $x\le 14$ is the number of points scored by the remaining 7 team members, we can simplify to obtain the Diophantine equation $x+15 = \frac{13}{28}t$, or $28x+28\cdot 15=13t$. Since $t$ is necessarily divisible by 28, let $t=28u$ where $u \ge 0$ and divide by 28 to obtain $x + 15 = 13u$. Then, it is easy to see $u=2$ ($t=56$) is the only candidate remaining, giving $x=\boxed{\textbf{(B)} 11}$. 

 -scrabbler94

 



\textbf{Solution 2}

We first start by setting the total number of points as $28$, since $\text{LCM}(4,7) = 28$. However, we see that this does not work since we surpass the number of points just with the information given ($28\cdot\frac{1}{4}+28\cdot\frac{2}{7} + 15 = 30$ $(> 28)$ ). Next, we can see that the total number of points scored is $56$ as, if it is more than or equal to $84$, at least one of the others will score more than 2 points. With this, we have that Alexa, Brittany, and Chelsea score: $56\cdot\frac{1}{4}+56\cdot\frac{2}{7} + 15 = 45$, and thus, the other seven players would have scored a total of $56-45 = \boxed{\textbf{(B)} 11}$. (We see that this works since we could have $4$ of them score $2$ points, and the other $3$ of them score $1$ point.) 

 -aops5234 -Edited by \href{https://artofproblemsolving.com/wiki/index.php?title=User:Penguin\_Spellcaster}{Penguin\_Spellcaster}

 



\textbf{Solution 3 — Modular Arithmetic}

Adding together Alexa's and Brittany's fractions, we get $\frac{15}{28}$ as the fraction of the total number of points they scored together. However, this is just a ratio, so we can introduce a variable: $\frac{15x}{28x}$ where $x$ is the common ratio. Let $y$ and $z$ and $w$ be the number of people who scored 1, 2, and 0 points, respectively. Writing an equation, we have $\frac{13x}{28x} = 15 + y + 2z + 0w.$ We want all of our variables to be integers. Thus, we want $15 + y + 2z = 0 \pmod {13}.$ Simplifying, $y+2z = 11 \pmod {13}.$ The only possible value, as this integer sum has to be less than $7 \cdot 2 + 1 = 15,$ must be 11. Therefore, $y+2z = 11,$ and the answer is $\boxed{ \textbf{(B) 11}}$.

 - ab2024

 



\textbf{Solution 4: Answer choices}

We can rewrite the question as an algebraic equation: $\frac{1}{4} T + \frac{2}{7} T + 15 + y = x$, where $x$ represents the total amount of points and $y$ the amount of points the $7$ other players scored. From there, we add the two fractions to get $\frac{15}{28} x + 15 + y = x$. Subtracting $\frac{15}{28} x$ from both sides, we get $\frac{13}{28} x = y + 15$. We multiply each side by $28$ to get rid of the denominator, in which we get $13x = 420 + 28y$. Now let’s think of this logically. This equation is telling us that if you add $420$ and $28$ times the amount of points scored by the extra $7$ players, you get $13$ times the amount of points total. And since we have to have a whole number of points total, this means that $420 + 28y$ must be divisible by $13$. Plugging in all the answer choices for $y$, we find that the only answer that makes $420 + 28y$ divisible by $13$ is $\boxed{ \textbf{(B) 11}}$.

 



\textbf{Solution 5 (modular congruences)}

As in the above, we can write the equation 
 \[\frac{1}{4}k+\frac{2}{7}k+15+y=k\]where $k$ is the total points, and $y$ is the rest of the students' total. This simplifies to the equation
 \[y+15=\frac{13}{28}k \implies\]
 \[28y+15 \cdot 28 = 13k\]Since $y,k$ are both integers, we have the modular congruence equation:
 \[28y+15 \equiv 0 \pmod{13} \implies\]
 \[2y \equiv 11 \pmod{13}\]Since the modular inverse of $2 \pmod{13}$ is $7$, this simplifies to:
 \[y \equiv 11 \cdot 7 \equiv 11 \pmod{13}\]The question asks for the least value of $y$. This is $11$, thus $\boxed{(B)11}$ 
 ~iretra

 



\textbf{Solution 6 (Because it doesn't have to be that complicated)}

We know from the problem that $\frac{1}{4} + \frac{2}{7}$ of the points are already taken. Then we add them together yielding, $\frac{15}{28}$, which means there are $\frac{13}{28}$ points left. Here is the magic part: Multiply it by 56, giving $26$. And now we can subtract Chelsea's $15$ points. That leaves us with $11$. Therefore, The answer is $\boxed{(B)11}$ 
P.S. Why do we multiply by $56$? Since the answer is a whole number we have to multiply by a multiple of $28$. We try different value and find $28$ is too small and $84$ is too big. 
~Fluffy1234

 


	\item (\textit{AMC 8 2019}) In triangle $ABC$, point $D$ divides side $\overline{AC}$ so that $AD:DC=1:2$. Let $E$ be the midpoint of $\overline{BD}$ and let $F$ be the point of intersection of line $BC$ and line $AE$. Given that the area of $\triangle ABC$ is $360$, what is the area of $\triangle EBF$?


 
                    \begin{figure}[H]
                        \centering
                        \includegraphics[width=0.5\textwidth]{Problem/AMC_8_AMC_2019_24_Problem_Figure_1.png}
                    \end{figure}
                    
$\textbf{(A) }24\qquad\textbf{(B) }30\qquad\textbf{(C) }32\qquad\textbf{(D) }36\qquad\textbf{(E) }40$

 \textbf{Answer: }B


	\item (\textit{AMC 8 2019}) Alice has 24 apples. In how many ways can she share them with Becky and
Chris so that each of the three people has at least two apples?


 $\textbf{(A) }105\qquad\textbf{(B) }114\qquad\textbf{(C) }190\qquad\textbf{(D) }210\qquad\textbf{(E) }380$

 \textbf{Answer: }C



\textbf{Solution 1 (Stars and Bars/Sticks and Stones)}

Note: This solution uses the non-negative version for stars and bars. A solution using the positive version of stars is similar (first removing an apple from each person instead of 2).

 This method uses the counting method of stars and bars (non-negative version). Since each person must have at least $2$ apples, we can remove $2*3$ apples from the total that need to be sorted. With the remaining $18$ apples, we can use stars and bars to determine the number of possibilities. Assume there are $18$ stars in a row, and $2$ bars, which will be placed to separate the stars into groups of $3$. In total, there are $18$ spaces for stars $+ 2$ spaces for bars, for a total of $20$ spaces. We can now do $20 \choose 2$. This is because if we choose distinct $2$ spots for the bars to be placed, each combo of $3$ groups will be different, and all apples will add up to $18$. We can also do this because the apples are indistinguishable. $20 \choose 2$ is $190$, therefore the answer is $\boxed{\textbf{(C) }190}$.

 



\textbf{Solution 2 (Answer Choices)}

Consider an unordered triple $(a,b,c)$ where $a+b+c=24$ and $a,b,c$ are not necessarily distinct. Then, we will either have $1$, $3$, or $6$ distinguishable ways to assign $a$, $b$, and $c$ to Alice, Becky, and Chris. Thus, our answer will be $x+3y+6z$ for some nonnegative integers $x,y,z$. Notice that we only have $1$ way to assign the numbers $a,b,c$ to Alice, Becky, and Chris when $a=b=c$. As this only happens $1$ way ($a=b=c=8$), our answer is $1+3y+6z$ for some $y,z$. Finally, notice that this implies the answer is $1$ mod $3$. The only answer choice that satisfies this is $\boxed{\textbf{(C) }190}$. 

 -BorealBear

 



\textbf{Solution 3}

Since each person needs to have at least two apples, we can simply give each person two, leaving $24 - 2\times3=18$ apples. For the remaining apples, if Alice is going to have $a$ apples, Becky is going to have $b$ apples, and Chris is going to have $c$ apples, we have indeterminate equation $a+b+c=18$. Currently, we can see that $0 \leq a\leq 18$ where $a$ is an integer, and when $a$ equals any number in the range, there will be $18-a+1=19-a$ sets of values for $b$ and $c$. Thus, there are $19 + 18 + 17 + \cdots + 1 = \boxed{\textbf{(C) }190}$ possible sets of values in total.

 


\end{enumerate}